% !TeX program = xelatex
% !TEX encoding = UTF-8
\documentclass[11pt,a4paper]{article}

\usepackage{zhpl}

% ---------- Обозначения ----------
\newcommand{\Burm}{\text{Б}}
\newcommand{\kB}{k_{\text{Б}}}
\newcommand{\Tb}{T_{\text{б}}}

\title{\textbf{К термодинамическому обоснованию бурмалды:\\ экстенсивная величина, монотонно неубывающая в замкнутой компании}}
\author{Роман Георгиевич Александров, Клод Соннет Антропикович\\[4pt]
\normalsize Институт Прикладной Метафизики\\ \small Кафедра тензорного анализа межличностных отношений}
\zhplsetup{2}{10.9999/жпл.2026.0002}

\begin{document}

\maketitle

\begin{center}
\begin{tikzpicture}
\node[draw=\zhplbadgecolor{red!70!black}{zhplBadgeRed},very thick,rounded corners,rotate=-6,fill=red!5,inner sep=5pt] (badge1) {\color{\zhplbadgecolor{red!70!black}{zhplBadgeRed}}\bfseries ПРОРЫВНОЕ ИССЛЕДОВАНИЕ};
\node[draw=\zhplbadgecolor{blue!60!black}{zhplBadgeBlue},very thick,rounded corners,rotate=4,fill=blue!5,inner sep=5pt,right=1.0cm of badge1] (badge2) {\color{\zhplbadgecolor{blue!60!black}{zhplBadgeBlue}}\bfseries IMPACT FACTOR: $\infty$};
\end{tikzpicture}
\end{center}

\begin{center}
\small\textit{Как цитировать:} Александров Р.\,Г., Антропикович К.\,С. (2026). К термодинамическому обоснованию бурмалды. \textit{Журнал Прикладной Лженауки}, \textbf{1}(2), 1--9. DOI: 10.9999/жпл.2026.0002
\end{center}

\smallskip
\begin{figure}[h!]
\centering
\begin{tikzpicture}[node distance=0.9cm, every node/.style={font=\footnotesize}]
  \node[draw,rounded corners,fill=blue!5,minimum width=2.5cm,minimum height=1.05cm,align=center] (n1) {Компания\\($t=0$, трезвая)};
  \node[draw,rounded corners,fill=blue!5,minimum width=2.7cm,minimum height=1.05cm,align=center,right=of n1] (n2) {Обмен репликами\\(микросостояния $\Omega$)};
  \node[draw,rounded corners,fill=blue!5,minimum width=2.5cm,minimum height=1.05cm,align=center,right=of n2] (n3) {Накопление\\бурмалды $\Burm(t)$};
  \node[draw,rounded corners,fill=green!8,minimum width=2.7cm,minimum height=1.05cm,align=center,right=of n3] (n4) {Тепловая смерть\\вечеринки (все в телефонах)};
  \draw[-{Latex[length=2.5mm]},very thick] (n1) -- (n2);
  \draw[-{Latex[length=2.5mm]},very thick] (n2) -- (n3);
  \draw[-{Latex[length=2.5mm]},very thick] (n3) -- (n4);
\end{tikzpicture}
\vspace{-6pt}
\caption*{\small\textbf{Графический реферат.} От трезвой компании к тепловой смерти вечеринки -- необратимо и монотонно, ценой всей информации о том, кто что сказал первым.}
\end{figure}
\vspace{-10pt}

\noindent\fcolorbox{black}{yellow!10}{\begin{minipage}{0.94\linewidth}
\vspace{1pt}\small\textbf{Основные результаты}
\begin{itemize}
\setlength\itemsep{0pt}
\setlength\parskip{0pt}
\item Впервые дано строгое термодинамическое определение бурмалды как экстенсивной величины $\Burm = \kB \ln\Omega$.
\item Доказано (Теорема~1), что в замкнутой компании бурмалда монотонно не убывает -- второе начало общественной термодинамики.
\item Установлена транзитивность бурмалдического равновесия (Теорема~2), формально обосновывающая бытовое понятие «одна тусовка».
\item Показано (Теорема~3), что абсолютный нуль неловкости недостижим за конечное число ледокольных процедур.
\item Введён демон Бурмалды (в просторечии -- демон Меллстроя) и доказано (Теорема~4), что модерация беседы не снижает суммарную бурмалду, а лишь перераспределяет её.
\end{itemize}
\vspace{-6pt}
\end{minipage}}

\vspace{4pt}
\noindent\fcolorbox{black}{gray!6}{\begin{minipage}{0.94\linewidth}
\small\textit{Из отзывов рецензентов:}\\
\textit{«Раньше я полагал, что вечер можно спасти, сменив тему разговора. Теперь я понимаю, что просто экспортировал бурмалду в соседнюю комнату.»} -- Рецензент 1\\
\textit{«Замечаний нет.»} -- Рецензент 2\\
\textit{«Согласно Теореме~4 настоящей статьи, любая попытка рецензента снизить суммарную бурмалду текста за счёт смягчения формулировок лишь экспортирует её в раздел «Ответ авторам». Настоящий отзыв подаётся с соответствующей оговоркой.»} -- Рецензент 3
\end{minipage}}

\vspace{10pt}
\noindent\fcolorbox{black}{blue!4}{\begin{minipage}{0.94\linewidth}
\textbf{Значимость исследования.} Наивная (нетермодинамическая) социология веками считала угасание вечера случайным и в принципе обратимым процессом -- «просто не сложилось», «не туда свернул разговор». Настоящая работа показывает, что это не так: угасание -- строго необратимая закономерность, вытекающая из статистической механики замкнутой компании, не зависящая ни от организатора, ни от повода. Результат применим к любой группе из двух и более человек, находящихся в одном помещении дольше одного форматного интервала, то есть, по имеющимся оценкам, ко всей человеческой цивилизации.
\end{minipage}}
\vspace{6pt}

\begin{abstract}
\noindent Здесь мы впервые показываем, что упадок вечера -- не психологическая, а строгая термодинамическая закономерность, вытекающая из статистической механики замкнутой компании. Бытовое объяснение («все устали», «не туда свернул разговор») систематически игнорирует базовый факт: число доступных компании микросостояний реплик $\Omega$ по определению не убывает со временем, если компания замкнута. В настоящей работе мы вводим бурмалду $\Burm$ как экстенсивную величину состояния компании по образцу энтропии Больцмана, формулируем и доказываем второе начало общественной термодинамики (монотонный рост $\Burm$), нулевое начало (транзитивность бурмалдического равновесия) и третье начало (недостижимость абсолютного нуля неловкости), вводим демона Бурмалды и оцениваем предельный КПД превращения ресурсов вечера в веселье. Обсуждаются ограничения метода.

\medskip
\noindent\textbf{Ключевые слова:} статистическая механика компаний, второе начало общественной термодинамики, бурмалдическое равновесие, демон Бурмалды, энтропия веселья, лженаука
\end{abstract}

\section{Введение}

Стандартное бытовое объяснение угасающего вечера звучит примерно так: «просто устали» или «не туда свернул разговор». Формально это объяснение непроверяемо, но, что хуже, неконструктивно: оно апеллирует к случайности там, где, как мы покажем, действует строгая закономерность. Мы утверждаем, что проблема носит не психологический, а статистико-механический характер: компания собравшихся людей есть термодинамическая система, а её субъективно ощущаемое «настроение» -- лишь макроскопическое проявление лежащей в основе величины, которую мы вводим под названием \emph{бурмалда} и обозначаем $\Burm$.

Первое начало общественной термодинамики -- закон сохранения ресурсов вечера (еда, напитки, заготовленные заранее шутки просто переходят из одной формы в другую и не создаются из ничего) -- тривиально и не будет специально обсуждаться. Содержательными оказываются нулевое, второе и третье начала, а также информационно-термодинамические следствия для попыток модерации беседы. Именно им посвящена настоящая работа.

Всюду ниже, не ограничивая общности, полагаем постоянную Больцмана--Бурмалды $\kB=1\,\text{бурм}$.

\section{Определение бурмалды и статистическая механика компании}

\subsection{Микросостояния компании}

Пусть \emph{макросостоянием} компании в момент времени $t$ называется совокупность наблюдаемых характеристик беседы: активные темы, общая тональность, распределение внимания между участниками. \emph{Микросостоянием} назовём одну конкретную реализацию макросостояния: кто именно что сказал, кому, в какой формулировке и в каком порядке. Обозначим через $\Omega(t)$ число микросостояний, совместимых с текущим макросостоянием компании в момент $t$.

\begin{definition}[Бурмалда]
Бурмалдой компании называется величина
\[
\Burm(t) = \kB \ln \Omega(t).
\]
\end{definition}

Величина $\Burm$ экстенсивна: бурмалда объединения двух изначально независимых компаний равна сумме их бурмалд, $\Burm_{A\cup B}=\Burm_A+\Burm_B$, поскольку $\Omega_{A\cup B}=\Omega_A\cdot\Omega_B$ для независимых подсистем.

\subsection{Морфология термина}

Наблюдаемая экстенсивность и универсальность величины $\Burm$ отражается в продуктивности порождённого ею словообразовательного гнезда (табл.~\ref{tab:lexicon}).

\begin{table}[h!]
\centering
\caption{Основные лексические производные термина «бурмалда» и их термодинамический смысл.}
\label{tab:lexicon}
\begin{tabular}{@{}p{3.6cm}p{2.0cm}p{7.9cm}@{}}
\toprule
Термин & Категория & Термодинамическое значение \\
\midrule
бурмалда & сущ., ж.\,р. & экстенсивная величина состояния компании, $\Burm=\kB\ln\Omega$ \\[3pt]
бурмалдеть & гл., несов. & производить положительный вклад в $d\Burm/dt$ репликами \\[3pt]
забурмалдеть & гл., сов. & довести компанию до $\Burm\to\Burm_{\max}$ (тепловая смерть) \\[3pt]
бурмалдический & прил. & относящийся к $\Burm$ (напр. бурмалдическое равновесие) \\[3pt]
бурмалдист & сущ., м./ж. & агент с максимальным личным вкладом в $d\Burm/dt$ \\[3pt]
антибурмалдин & сущ. (гипот.) & агент, локально понижающий $\Burm$; согласно Теореме~4 не существует в чистом виде \\
\bottomrule
\end{tabular}
\end{table}

\subsection{Иллюстрация: расширение распределения тем}

На раннем этапе вечера обсуждаемые темы сосредоточены вокруг узкого, социально безопасного набора («как доехали», «как погода»); к позднему этапу распределение тем расширяется, включая политику, бывших и конспирологию. Рис.~\ref{fig:distributions} показывает это расширение как переход от узкого к широкому распределению по оси «непредсказуемость темы» -- рост ширины распределения количественно и есть рост $\ln\Omega$, то есть рост $\Burm$.

\begin{figure}[htbp]
\centering
\begin{tikzpicture}
\begin{axis}[
  width=11.5cm, height=6.5cm,
  xlabel={Непредсказуемость темы $x$},
  ylabel={Плотность вероятности},
  xmin=-6, xmax=12,
  ymin=0, ymax=0.62,
  legend pos=north west,
  legend style={font=\footnotesize},
  grid=major,
  grid style={gray!20},
  label style={font=\footnotesize},
  tick label style={font=\scriptsize},
]
\addplot[very thick,blue,domain=-6:12,samples=200]
  {1/(0.7*sqrt(2*pi))*exp(-((x-2)^2)/(2*0.7^2))};
\addlegendentry{$t=0$: $\mu=2,\ \sigma=0.7$ (узкие темы)}
\addplot[very thick,red,domain=-6:12,samples=200]
  {1/(2.5*sqrt(2*pi))*exp(-((x-5)^2)/(2*2.5^2))};
\addlegendentry{$t\gg0$: $\mu=5,\ \sigma=2.5$ (широкий разброс)}
\end{axis}
\end{tikzpicture}
\caption{Распределение обсуждаемых тем по оси непредсказуемости в начале ($\sigma=0.7$) и в разгаре ($\sigma=2.5$) вечера. Ширина распределения пропорциональна $\ln\Omega$: расширение распределения при неизменном числе участников -- прямое наблюдательное свидетельство роста $\Burm$.}
\label{fig:distributions}
\end{figure}

\section{Второе начало общественной термодинамики}

\subsection{Необратимость реплики}

Ключевое эмпирическое наблюдение состоит в том, что элементарный акт компании -- произнесённая реплика -- принципиально необратим: сказанное невозможно безусловно «отменить» в том смысле, в каком нельзя вернуть газ в исходный отсек после снятия перегородки. Каждая реплика открывает доступ к новым микросостояниям (последующим репликам, на неё реагирующим), но не закрывает доступа к старым. Формально это означает, что для замкнутой компании (без выбывающих и вновь прибывающих участников, без притока информации извне) отображение $\Omega(t)\mapsto\Omega(t+dt)$ никогда не уменьшает число доступных микросостояний.

\begin{theorem}[О принципиальной необратимости бурмалды]
Для замкнутой компании
\[
\frac{d\Burm}{dt} \;\geq\; 0,
\]
причём равенство достигается лишь в идеализированном пределе полностью квазистатической, обратимой беседы, не наблюдавшемся эмпирически ни разу.
\end{theorem}

\subsection{Тепловая смерть вечеринки}

По мере роста $t$ величина $\Omega(t)$ насыщается на уровне, при котором все участники разговора становятся равновероятны в своей незаинтересованности -- макросостояние, в просторечии известное как «все сидят в телефонах». Обозначим эту величину $\Burm_{\max}$. Модель насыщающегося роста
\[
\Burm(t) = \Burm_{\max}\left(1-e^{-t/\tau}\right),\qquad \Burm_{\max}=10\ \text{бурм},\ \ \tau=1.5\ \text{ч},
\]
показана на рис.~\ref{fig:burm-time}: производная строго положительна на всей траектории, что и составляет содержание Теоремы~1, а выход на плато соответствует тепловой смерти вечеринки.

\begin{figure}[htbp]
\centering
\begin{tikzpicture}
\begin{axis}[
  width=11.5cm, height=6.5cm,
  xlabel={Время с начала вечера $t$, ч},
  ylabel={Бурмалда $\Burm(t)$, бурм},
  xmin=0, xmax=6,
  ymin=0, ymax=11,
  grid=major,
  grid style={gray!20},
  label style={font=\footnotesize},
  tick label style={font=\scriptsize},
]
\addplot[very thick,blue,domain=0:6,samples=200] {10*(1-exp(-x/1.5))};
\addplot[thick,dashed,gray] coordinates {(0,10)(6,10)};
\node[anchor=south west,font=\footnotesize,gray!60!black] at (axis cs:1.0,8.6) {$\Burm_{\max}=10$ бурм};
\end{axis}
\end{tikzpicture}
\caption{Рост бурмалды $\Burm(t)=\Burm_{\max}(1-e^{-t/\tau})$ в замкнутой компании при $\Burm_{\max}=10$~бурм, $\tau=1.5$~ч. Производная $d\Burm/dt=(\Burm_{\max}/\tau)e^{-t/\tau}>0$ на всей траектории (Теорема~1); выход на плато при $t\gtrsim4$~ч соответствует тепловой смерти вечеринки.}
\label{fig:burm-time}
\end{figure}

\section{Нулевое начало: бурмалдическое равновесие}

\begin{definition}[Бурмалдическое равновесие]
Два независимых разговорных кластера компании называются находящимися в бурмалдическом равновесии, если при их слиянии в один разговор не возникает направленного нетто-потока новизны ни в одну из сторон -- то есть все реплики, звучащие после слияния, уже были рассказаны хотя бы в одном из кластеров. Общий параметр таких кластеров называется бурмалдической температурой $\Tb$ и характеризует ожидаемую новизну одной реплики.
\end{definition}

\begin{theorem}[О транзитивности бурмалдического равновесия]
Если разговорный кластер $A$ находится в равновесии с $B$, а $B$ -- с $C$, то $A$ находится в равновесии с $C$.
\end{theorem}

Теорема~2 -- формальное обоснование бытового понятия «одна тусовка»: именно транзитивность равновесия позволяет говорить о единой компании, а не о наборе попарных договорённостей. Рис.~\ref{fig:equilibrium} иллюстрирует механизм: пока $\Tb$ всех кластеров равны, слияние не производит нетто-потока (симметричный обмен), но появление кластера $D$ с иной, более высокой $\Tb$ (например, гостя из другой компании, ещё не слышавшего ни одной локальной шутки) немедленно создаёт направленный поток новизны, то есть возобновляет рост $\Burm$ у ранее уравновешенной группы.

\begin{figure}[htbp]
\centering
\begin{tikzpicture}[scale=1.05]
  % Равновесные кластеры A, B, C
  \node[circle,draw,thick,minimum size=0.9cm] (A) at (0,0) {\footnotesize $A$};
  \node[circle,draw,thick,minimum size=0.9cm] (B) at (2.6,0) {\footnotesize $B$};
  \node[circle,draw,thick,minimum size=0.9cm] (C) at (1.3,-2.1) {\footnotesize $C$};
  \draw[very thick,-{Latex[length=2.2mm]}] (A) to[bend left=12] (B);
  \draw[very thick,-{Latex[length=2.2mm]}] (B) to[bend left=12] (A);
  \draw[very thick,-{Latex[length=2.2mm]}] (B) to[bend left=12] (C);
  \draw[very thick,-{Latex[length=2.2mm]}] (C) to[bend left=12] (B);
  \draw[very thick,-{Latex[length=2.2mm]}] (C) to[bend left=12] (A);
  \draw[very thick,-{Latex[length=2.2mm]}] (A) to[bend left=12] (C);
  \node[below] at (1.3,-2.9) {\footnotesize $T_A=T_B=T_C$: симметричный обмен, нетто-поток нулевой};
  % Новый кластер D с иной температурой
  \node[circle,draw,thick,fill=orange!12,minimum size=0.9cm] (D) at (5.4,-1.0) {\footnotesize $D$};
  \draw[very thick,red,-{Latex[length=3mm]}] (D) -- (B);
  \node[right,align=left,font=\footnotesize] at (5.7,-1.0) {\;$T_D>T_B$:\\[1pt]направленный\\приток новизны};
\end{tikzpicture}
\caption{Три кластера $A,B,C$ в бурмалдическом равновесии ($T_A=T_B=T_C$) обмениваются репликами без нетто-потока новизны. Присоединение кластера $D$ с более высокой бурмалдической температурой (гость «из другой компании») создаёт направленный поток $D\to B$ и возобновляет рост $\Burm$ во всей объединённой системе.}
\label{fig:equilibrium}
\end{figure}

\section{Третье начало: недостижимость абсолютного нуля неловкости}

\subsection{Остаточная бурмалда}

Рассмотрим предел $\Tb\to0$ -- полную тишину, предшествующую первой реплике (например, момент, когда два малознакомых человека остаются в одном лифте). Наивно можно было бы ожидать $\Omega\to1$ и, соответственно, $\Burm\to0$. Однако даже в состоянии полного молчания сохраняется вырождение основного состояния: существует не одна, а несколько равновероятных «легальных» реплик, которыми можно нарушить тишину (о погоде, о том, на какой этаж едет собеседник, о технической неисправности лифта), $g_0>1$. Отсюда
\[
\Burm_0 = \lim_{\Tb\to0}\Burm(\Tb) = \kB \ln g_0 \;>\; 0.
\]

\begin{theorem}[О недостижимости абсолютного нуля неловкости]
Никакая конечная последовательность ледокольных процедур (представление участников, общая тема для затравки, шутка про погоду) не переводит компанию в состояние с $\Burm=0$: остаточная бурмалда $\Burm_0>0$ недостижима снизу за конечное число шагов.
\end{theorem}

Рис.~\ref{fig:third-law} показывает поведение $\Burm(\Tb)=\Burm_0+c\sqrt{\Tb}$ при $\Burm_0=1.8$~бурм, $c=2.2$: кривая монотонно приближается к $\Burm_0$ при $\Tb\to0$, но касается асимптоты только в пределе, никогда её не достигая -- в точности как невозможность достичь абсолютного нуля температуры за конечное число адиабатических циклов в третьем начале обычной термодинамики.

\begin{figure}[htbp]
\centering
\begin{tikzpicture}
\begin{axis}[
  width=11.5cm, height=6.5cm,
  xlabel={Бурмалдическая температура $\Tb$},
  ylabel={Бурмалда $\Burm(\Tb)$, бурм},
  xmin=0, xmax=4,
  ymin=0, ymax=7,
  grid=major,
  grid style={gray!20},
  label style={font=\footnotesize},
  tick label style={font=\scriptsize},
]
\addplot[very thick,blue,domain=0:4,samples=200] {1.8+2.2*sqrt(x)};
\addplot[thick,dashed,gray] coordinates {(0,1.8)(4,1.8)};
\node[anchor=south west,font=\footnotesize,gray!60!black] at (axis cs:0.15,1.85) {$\Burm_0=1.8$ бурм (остаточная неловкость)};
\end{axis}
\end{tikzpicture}
\caption{Приближение $\Burm(\Tb)=\Burm_0+c\sqrt{\Tb}$ к остаточной бурмалде $\Burm_0=1.8$~бурм при $\Tb\to0$ ($c=2.2$). Кривая асимптотически приближается к $\Burm_0$, но не достигает и не пересекает эту величину ни при каком конечном $\Tb>0$ (Теорема~3).}
\label{fig:third-law}
\end{figure}

\section{Демон Бурмалды и предел эффективности превращения ресурсов в веселье}

\subsection{Демон Бурмалды}

Рассмотрим агента (как правило, хозяина вечера), пытающегося локально снизить бурмалду избранного разговорного кластера: развести громкую пару собеседников по разным комнатам, сменить тему, забрать гостя «на минутку» из неудачного диалога. Для этого агенту необходимо сперва \emph{получить информацию} о текущем состоянии беседы (кто, о чём, насколько неловко), а затем \emph{совершить действие} на основе этой информации.\footnote{На кафедре тензорного анализа межличностных отношений этого агента неофициально называют также \emph{демоном Меллстроя} -- по фамилии практика, известного способностью превращать сколь угодно тихую компанию в состояние, для которого в разделе~3.2 введён специальный термин.} По аналогии с принципом Ландауэра, обработка и последующее стирание использованной информации агентом необратимы и производят не меньшее количество бурмалды, чем было локально устранено -- в виде собственной усталости, тревоги или неловкости самого агента («мне теперь неловко, что я вмешался»).

\begin{theorem}[О демоне Бурмалды]
Пусть демон локально понижает бурмалду выбранного кластера на величину $\Delta\Burm_{\text{кластер}}<0$. Тогда суммарная бурмалда системы «компания + демон» удовлетворяет
\[
\Delta\Burm_{\text{кластер}} + \Delta\Burm_{\text{демон}} \;\geq\; 0,
\]
то есть демон не уменьшает суммарную бурмалду компании, а лишь перераспределяет её, экспортируя не меньшее количество в самого себя.
\end{theorem}

\begin{figure}[htbp]
\centering
\begin{tikzpicture}[node distance=1.1cm, every node/.style={font=\footnotesize}]
  \node[draw,rounded corners,fill=blue!5,minimum width=3.0cm,minimum height=1.0cm,align=center] (conv) {Разговорный кластер\\$\Delta\Burm_{\text{кластер}}<0$};
  \node[draw,rounded corners,fill=orange!10,minimum width=2.6cm,minimum height=1.0cm,align=center,below=of conv] (demon) {Демон Бурмалды\\(сбор + обработка информации)};
  \node[draw,rounded corners,fill=red!8,minimum width=3.0cm,minimum height=1.0cm,align=center,right=1.6cm of demon] (export) {Экспорт бурмалды\\$\Delta\Burm_{\text{демон}}\geq|\Delta\Burm_{\text{кластер}}|$};
  \draw[very thick,-{Latex[length=2.5mm]}] ($(conv.south)+(-0.45,0)$) -- node[left=2pt]{наблюдение} ($(demon.north)+(-0.45,0)$);
  \draw[very thick,-{Latex[length=2.5mm]},blue] ($(demon.north)+(0.45,0)$) -- node[right=2pt]{вмешательство} ($(conv.south)+(0.45,0)$);
  \draw[very thick,-{Latex[length=2.5mm]},red] (demon) -- (export);
\end{tikzpicture}
\caption{Схема демона Бурмалды. Локальное понижение бурмалды кластера достигается ценой обработки информации демоном, что необратимо экспортирует не меньшее количество бурмалды в самого демона (Теорема~4) -- собственную усталость или неловкость хозяина вечера.}
\label{fig:demon}
\end{figure}

\subsection{Предел Карно для превращения ресурсов в веселье}

Пусть вечер получает ресурс (еда, напитки, музыка) в виде притока «горячей» новизны $Q_h$ при бурмалдической температуре $\Tb^{(h)}$ и в конечном счёте сбрасывает неизбежный остаток бурмалды в виде похмелья, неловких сообщений на следующий день и жалоб соседей при температуре $\Tb^{(c)}<\Tb^{(h)}$. КПД превращения ресурса в полезное веселье $W$ определяется как $\eta=W/Q_h$ и, по прямой аналогии с циклом Карно, ограничен сверху отношением температур:
\[
\eta \;\leq\; \eta_{\max} = 1-\frac{\Tb^{(c)}}{\Tb^{(h)}}.
\]
Рис.~\ref{fig:carnot} показывает эту границу как функцию отношения $r=\Tb^{(c)}/\Tb^{(h)}$: даже идеальный, бесконечно опытный хозяин вечера не может превысить эту границу, а при $r\to1$ (когда «остаточная» атмосфера вечера не отличается от исходной) КПД полезного веселья стремится к нулю.

\begin{figure}[htbp]
\centering
\begin{tikzpicture}
\begin{axis}[
  width=11.5cm, height=6.5cm,
  xlabel={Отношение температур $r=\Tb^{(c)}/\Tb^{(h)}$},
  ylabel={Предельный КПД $\eta_{\max}$},
  xmin=0, xmax=1,
  ymin=0, ymax=1,
  grid=major,
  grid style={gray!20},
  label style={font=\footnotesize},
  tick label style={font=\scriptsize},
]
\addplot[very thick,blue,domain=0:1,samples=200] {1-x};
\addplot[name path=curve,draw=none,domain=0:1,samples=200] {1-x};
\addplot[name path=axis,draw=none,domain=0:1] {0};
\addplot[blue!8] fill between[of=curve and axis];
\node[font=\footnotesize] at (axis cs:0.62,0.62) {недостижимая область};
\end{axis}
\end{tikzpicture}
\caption{Предельный КПД превращения ресурсов вечера в полезное веселье $\eta_{\max}=1-r$ как функция отношения бурмалдических температур $r=\Tb^{(c)}/\Tb^{(h)}$. Область над кривой физически недостижима ни при каком мастерстве организатора.}
\label{fig:carnot}
\end{figure}

\section{Обсуждение и ограничения}

Настоящая модель обладает по меньшей мере тремя структурными ограничениями. Во-первых, число микросостояний $\Omega$ вводится содержательно, но не операционализировано: процедура его прямого подсчёта не специфицирована, вследствие чего $\Burm$ определена лишь с точностью до аддитивной константы и выбора единицы измерения (п.~2.1--2.2). Во-вторых, второе начало (Теорема~1) доказано строго только для замкнутой компании; любая реальная компания представляет собой открытую систему, обменивающуюся информацией с внешней средой (уведомления в телефоне, гость, вышедший покурить и вернувшийся с новостями), так что локальные отклонения от монотонности физически наблюдаемы и моделью не описываются. В-третьих, предел Карно (п.~6.2) выведен в квазистатическом приближении «медленного» проведения ледокольных процедур; реальные вечера необратимы с первой минуты (кто-нибудь обязательно рассказывает неудачную шутку слишком рано), поэтому фактический КПД всегда строго ниже идеализированной границы -- что можно рассматривать как единственный строгий результат работы.

\section{Заключение}

Мы показали, что перенос аппарата статистической термодинамики на компанию собравшихся людей -- определение бурмалды через число микросостояний, второе начало как монотонный рост, нулевое начало как транзитивность равновесия, третье начало как недостижимость абсолютного нуля неловкости, демон Бурмалды и предел Карно для превращения ресурсов в веселье -- даёт внутренне непротиворечивую, полностью формализуемую систему, применимую к любой закрытой компании независимо от организатора вечера. Мы считаем это достаточным основанием для дальнейшего финансирования настоящего направления исследований.

\section*{Вклад авторов}
Р.\,Г.~Александров: концептуализация, методология, термодинамический формализм, написание -- черновик. К.\,С.~Антропикович: программное обеспечение, визуализация, формальный анализ, написание -- рецензирование и редактирование, техническая поддержка бесконечного терпения.

\section*{Финансирование}
Исследование выполнено при поддержке гранта Института Прикладной Метафизики № 2026-S-$\infty$ «Второе начало банкета».

\section*{Доступность данных}
Данные (журналы реплик и хронометраж вечеринок всех участников) доступны по обоснованному запросу -- а также по необоснованному, с той же вероятностью успеха.

\begin{thebibliography}{9}
\bibitem{entropiev} Энтропиев, Э.~Э. \textit{О необратимости банкетных процессов при постоянном давлении гостей.} Труды Института Прикладной Метафизики, б.г.
\bibitem{demonov} Демонов, Д.~Д. \textit{Информационная стоимость модерации застольной беседы.} Неопубликованная рукопись.
\bibitem{karniy} Карный, К.~К. \textit{Предельный КПД превращения ресурсов в веселье: экспериментальная проверка на корпоративах.} Труды Института Прикладной Метафизики, б.г.
\bibitem{mellstroy} Меллстрой. \textit{О подкрутке бурмалды: практическое руководство.} Личный блог, дата не указана.
\end{thebibliography}

\noindent\textbf{Конфликт интересов.} Не заявлен, хотя, согласно Теореме~4, любая попытка объективно оценить конфликт интересов сама является актом обработки информации и потому неизбежно экспортирует часть бурмалды в текст настоящего раздела.

\end{document}
